
Benchmark contamination — the overlap between NLP test sets and training corpora — is widely recognized as a potential source of evaluation score inflation, yet existing analyses treat contamination as a single scalar per benchmark rather than a structured phenomenon that varies across sub-tasks and corpora. This paper presents a systematic multi-corpus, multi-sub-task 13-gram contamination analysis: a 59-sub-task × 3-corpus matrix of 13-gram containment rates for MMLU, HellaSwag, and BIG-Bench Hard against The Pile v1, C4 en.noclean, and RedPajama-v1. The central finding is that contamination rates vary by approximately 40× across benchmark sub-tasks within a single corpus (Kruskal-Wallis H=590.82, p=2.73×10⁻⁸⁹), with professional_medicine at 17.3\% and high_school_mathematics at 0.4\% in The Pile v1. Across corpora, a statistically significant ordering is observed (Kruskal-Wallis H=17.51, p=1.58×10⁻⁴): The Pile (mean 6.53\%) > RedPajama (5.75\%) > C4 (4.05\%). The Pile and RedPajama are statistically indistinguishable (Dunn post-hoc p=0.810), consistent with their shared CommonCrawl ancestry, while C4's quality filtering is associated with a 38\% lower mean contamination rate relative to The Pile (assuming literature-calibrated scaling factors applied to 10\% corpus samples; directional ordering robust to ρ>0.995 sensitivity analysis). A domain-stratified analysis finds that academic MMLU sub-tasks show approximately 1.9× higher contamination than commonsense benchmarks across all three corpora (Cohen's d=0.85), but this directional pattern did not achieve Mann-Whitney statistical significance due to insufficient commonsense sub-task granularity (n=2). A fourth planned experiment on metric consistency between 13-gram containment and Jaccard similarity was not executed. These results indicate that corpus source composition is associated with predictable, structured contamination profiles across benchmark sub-tasks.

